\section{神经网络的本质：函数逼近器}

\subsection{Universal Approximation Theorem}

要理解深度学习的威力与局限，一个非常重要的出发点是 1989 年 Hornik 等人提出的 \textbf{Universal Approximation Theorem}（万能逼近定理）：

\begin{importantbox}{Universal Approximation Theorem}
一个具有单隐层的前馈神经网络，在隐藏单元数量足够多的情况下，能够以任意精度逼近任何连续函数。
\end{importantbox}

$$
f(x) \approx \hat{f}(x; W, b) = \sigma(W_2 \cdot \sigma(W_1 x + b_1) + b_2)
$$

其中：
\begin{itemize}
    \item $f(x)$：目标连续函数
    \item $\hat{f}(x; W, b)$：单隐层神经网络的输出
    \item $\sigma$：非线性激活函数
    \item $W_1, W_2, b_1, b_2$：可学习的权重和偏置参数
\end{itemize}

\begin{figure}[H]
\centering
\includegraphics[width=0.85\textwidth]{slides-images/slide-12.jpg}
\caption{Universal Approximation Theorem 示意图\protect\footnotemark}
\end{figure}
\footnotetext{Slides 第12页。引用 Hornik+ Neural Networks 1989。}

\subsection{定理的局限}

虽然该定理在理论上非常强大，但存在两个关键的 caveat：

\begin{figure}[H]
\centering
\includegraphics[width=0.85\textwidth]{slides-images/slide-13.jpg}
\caption{Universal Approximation Theorem 的两个重要注意事项\protect\footnotemark}
\end{figure}
\footnotetext{Slides 第13页。}

\begin{warningbox}{万能逼近定理的两大注意事项}
\begin{enumerate}
    \item \textbf{隐藏单元数量可能不可行地大}：定理保证存在这样一个网络，但没有给出所需神经元数量的上界——对于复杂函数，可能需要天文数字级别的参数。
    \item \textbf{不保证泛化能力}：定理只保证逼近能力，不保证学到的模型能在未见数据上表现良好。拟合训练数据不等于理解数据的底层规律。
\end{enumerate}
\end{warningbox}

此外，定理并未提供如何找到最优权重的方法论。梯度下降虽然是一种有效的优化策略，但在高维非凸损失曲面上，优化本身就是一个高度非平凡的问题。

\subsection{AI 发展的历史视角}

\begin{figure}[H]
\centering
\includegraphics[width=0.85\textwidth]{slides-images/slide-14.jpg}
\caption{AI ``Hype'' 的历史脉络\protect\footnotemark}
\end{figure}
\footnotetext{Slides 第14页。}

AI 的发展史经历了典型的 hype cycle：从 1956 年达特茅斯会议的诞生，到 1974--1980 年的第一次 AI 寒冬，再到 1987--1993 年的第二次寒冬，最终在深度学习时代迎来爆发式增长。理解这段历史有助于我们对当前的 AI 热潮保持清醒的判断——既要拥抱技术进步，也要警惕过度炒作。

\subsection{本章小结}

神经网络从本质上说是函数逼近器。万能逼近定理保证了理论上的表达能力，但在实际应用中，隐藏层大小、优化难度和泛化能力都是必须面对的现实挑战。这一认识框架为下一节讨论深度学习的具体局限性奠定了基础。

%% =====================================================
